
\exercise
Give an example of two commuting operators $S, T$ on $\F4$ such that there is a subspace of $\F4$ that is invariant under $S$ but not under $T$ and there is a subspace of $\F4$ that is invariant under $T$ but not under $S$.


\exercise
Suppose $\E$ is a subset of $\L(V)$ and every element of $\E$ is diagonalizable. Prove that there exists a basis of $V$ with respect to which every element of $\E$ has a diagonal matrix if and only if every pair of elements of $\E$ commutes.%
 \label{3Diag}\index{commuting operators|(}
\exercisecomment{This exercise extends \ref{8SimDiag}, which considers the case in which\/ $\E$ contains only two elements. For this exercise,\/ $\E$ may contain any number of elements, and\/ $\E$ may even be an infinite set.}


\exercise
Suppose $S, T \in \L(V)$ are such that $ST = TS$. Suppose $p \in \P(\F{})$.
\begin{subtheorem}
\item
Prove that $\ker p(S)$ is invariant under $T$.
\item
Prove that $\ran p(S)$ is invariant under $T$.
\end{subtheorem}
\exercisecomment{See \ref{130} for the special case $S = T$.}


\exercise
Prove or give a counterexample: If $A$ is a diagonal matrix and $B$ is an upper-triangular matrix of the same size as $A$, then $A$ and $B$ commute.


\exercise
Prove that a pair of operators on a finite-dimensional vector space commute if and only if their dual operators commute.
\exercisecomment{See \ref{DualMap} for the definition of the dual of an operator.}\label{5dualcom}


\exercise
Suppose that $V$ is a nonzero finite-dimensional complex vector space and $S, T \in \L(V)$ commute. Prove that there exist $\al, \la \in \C{}$ such that
\[
\ran (S - \al I) + \ran (T - \la I) \ne V\1.
\]
\exercisedisplayend


\exercise
Suppose $V$ is a complex vector space, $S \in \L(V)$ is diagonalizable, and $T \in \L(V)$ commutes with $S$. Prove that there is a basis of $V$ such that $S$ has a diagonal matrix with respect to this basis and $T$ has an upper-triangular matrix with respect to this basis.


\exercise
Suppose $m = 3$ in Example \ref{PD} and $D_x, D_y$ are the commuting partial differentiation operators on $\P\!_3\3\paren[\big]{\R2}$ from that example. Find a basis of $\P\!_3\3\paren[\big]{\R2}$ with respect to which $D_x$ and $D_y$ each have an upper-triangular matrix.


\exercise
Suppose $V$ is a finite-dimensional nonzero complex vector space. Suppose that $\E \subset \L(V)$ is such that $S$ and $T$ commute for all $S, T \in \E$. \label{4ManyOps}
\begin{subtheorem}
\item
Prove that there is a vector in $V$ that is an eigenvector for every element of $\E$.
\item
Prove that there is a basis of $V$ with respect to which every element of $\E$ has an upper-triangular matrix.
\end{subtheorem}
\exercisecomment{This exercise extends \ref{2common} and \ref{2commute}, which consider the case in which\/ $\E$ contains only two elements. For this exercise,\/ $\E$ may contain any number of elements, and\/ $\E$ may even be an infinite set.}%


\exercise
Give an example of two commuting operators $S, T$ on a finite-dimensional real vector space such that $S+T$ has an eigenvalue that does not equal an eigenvalue of $S$ plus an eigenvalue of $T$ and $ST$ has an eigenvalue that does not equal an eigenvalue of $S$ times an eigenvalue of $T\3$.\index{commuting operators|)}
\exercisecomment{This exercise shows that \ref{2EigenSum} does not hold on real vector spaces.}


